Define the notion of a monetary risk measure. What is the corresponding accep-
tance set? Explain the important notion of capital requirements.

Monetary Risk Measure
A mapping \(\rho : \mathcal{X} \to \mathbb{R}\) is called a monetary risk measure if it satisfies the following properties:  
1. Inverse monotonicity: If \(\mathcal{X}(\omega) \leq \mathcal{Y}(\omega)\) for all \(\omega \in \Omega\), then 
   \(\rho(\mathcal{X}) \geq \rho(\mathcal{Y})\).  
2. Cash invariance: For any \(m \in \mathbb{R}\), \(\rho(\mathcal{X} + m) = \rho(\mathcal{X}) - m\).

Interpretation:  
1. Property 1 states that the risk of a position \(\mathcal{Y}\) is smaller than the risk of a position \(\mathcal{X}\), 
   if the future value of \(\mathcal{Y}\) is at least as large as \(\mathcal{X}\) for any scenario.  
2. Property 2 states that risk is measured on a monetary scale: if \(m\) are added to \(\mathcal{X}\), 
   then the risk of \(\mathcal{X}\) is exactly reduced by this amount.

Acceptability:  
- A financial position \(X \in \mathcal{X}\) is said to be acceptable with respect to a given monetary risk measure \(\rho\) if \(\rho(X) \leq 0\).  
- The set \(\mathcal{A} = \mathcal{A}_{\rho}\) of all acceptable positions is called the acceptance set:  
  \( \mathcal{A} := \{X \in \mathcal{X} \mid \rho(X) \leq 0\} \)

The acceptance set is non-empty and satisfies the following two conditions:  
1. \(\inf \{m \in \mathbb{R} \mid m \in \mathcal{A}\} > -\infty\)  
2. \(X \in \mathcal{A}, Y \in \mathcal{X}, Y \geq X \Rightarrow Y \in \mathcal{A}\)

The monetary risk measure \(\rho\) can be recovered from its acceptance set  
\(\rho(X) = \inf \{m \in \mathbb{R} \mid X + m \in \mathcal{A}\}\),

i.e., \(\rho\) is indeed a capital requirement.


Define and discuss the notions convex risk measure, coherent risk measure, sub-
additive risk measure. How do these properties correspond to properties of the
acceptance set?

Properties of Acceptance Sets

Let \(\rho\) be a monetary risk measure with acceptance set \(\mathcal{A}\).

- \(\rho\) is a convex risk measure if and only if \(\mathcal{A}\) is convex.
- \(\rho\) is positively homogeneous if and only if \(\mathcal{A}\) is a cone.
- \(\rho\) is coherent if and only if \(\mathcal{A}\) is a convex cone.

Acceptance sets as a starting point for regulation:

- Acceptance sets divide the universe of all positions into those that are acceptable and those that are unacceptable.
- They correspond to risk measures which possess an operational interpretation as capital requirements.
- Thus, acceptance sets are a good starting point to formulate the normative standards of a regulator.

What is the difference of measures of dispersion on the one hand and downside
risk measures on the other? Give examples

Measures of dispersion:
Variance:  
\(
\text{Var}(X) = \mathbb{E}[(X - \mathbb{E}[X])^2]
\)
- The variance quantifies the expected quadratic deviation between a random variable and its expectation.
- In the sense of the two-sided definition of risk deviations of \(X\) from \(\mathbb{E}[X]\) in both directions are taken into account.

Standard deviation (volatility): \(\sigma(X)\) quantifies risk in the same unit as \(X\) itself.
\(
\sigma(X) = \sqrt{\text{Var}(X)} = \sqrt{\mathbb{E}[(X - \mathbb{E}[X])^2]}
\)

One-sided standard deviation (semi-standard deviation):
\(
\sigma_{+}(X) = \sqrt{\mathbb{E}[(\mathbb{E}[X] - X)_{+}^2]}
\)
\(\sigma_{+}(L) = \sqrt{\mathbb{E}[(L - \mathbb{E}[L])_{+}^2]}\)
In contrast to the variance or the standard deviation, the one-side standard deviation only accounts for shortfalls \(X < \mathbb{E}[X]\) 
in case of a financial position \(X\) and for exceedances \(L > \mathbb{E}[L]\) in case of loss variables respectively.


Random variables with a higher expectation typically show a higher variance in comparison to
random variables with smaller expectation.
Conclusion: The interpretation of the variance or standard deviation is only meaningful in standardized terms.
Coefficient of variation:
\(
\text{cv}(X) = \frac{\sigma(X)}{\mathbb{E}[X]}
\)

Shortfall measures quantify the risk of exceeding a given loss threshold \( a \).

Example 4.1.3 (Upper and Lower Partial Moments)

- The upper partial moments refer to loss variables and weight the exceedances by a power function:
\[
\text{UPM}_{k,a}(L) :=
\begin{cases} 
\mathbb{E}[(L - a)^+]^k & \text{for } k > 0, \\
\mathbb{P}[L \ge a] & \text{for } k = 0.
\end{cases}
\]

- In the same manner, the lower partial moments are defined with respect to financial positions:
\[
\text{LPM}_{k,a}(X) :=
\begin{cases} 
\mathbb{E}[(a - X)^+]^k & \text{for } k > 0, \\
\mathbb{P}[X < a] & \text{for } k = 0.
\end{cases}
\]

Special cases:
- \( k = 0 \): Shortfall probability with respect to the loss threshold \( a \),
- \( k = 1 \): expected exceedance of the loss threshold \( a \),
- \( k = 2 \): shortfall variance.


Explain scenario based risk measures

Scenario-Based Risk Measures for Stress Testing

Philosophy: Risk is measured as maximum (weighted) loss under certain scenarios.

- Scenarios typically correspond to risk factor changes, e.g., 
  yield curve shift by 50 basis points, increase of credit spreads, increase of mortality rates etc.
- In practice, it is a common approach to consider cross scenarios, i.e., simultaneous changes of risk factors.

Notation:
- Set of hypothetical scenarios \( S = \{s_1, \ldots, s_n\} \) with corresponding weights \((w_1, \ldots, w_n)\)
- \( L(s_i) \) portfolio loss for scenario \( s_i \), \( i = 1, \dots, n\)

Definition: Scenario based risk measure

\( \rho_S \coloneqq \text{max}_{i} \{w_i L(s_i)\}\)

Assume that \( L(0) = 0 \) (normalization) and \( w_i \in [0, 1] \)
Consider \( S \sim \mathbb{P}_i \) for the probability measure \( \mathbb{P} \) with \(\mathbb{P}_i[S = si ] = wi\) and 
\(\mathbb{P}_i[S = 0] = 1 - wi\)
\( w_iL(s_i) = w_iL(s_i) + (1 − w_i)L(0) = E_{\mathbb{P}_i}[L(S)] \).
Then, \( \rho_S = \text{max}_{\mathbb{P} \in \{\mathbb{P}_1, \dots, \mathbb{P}_n\} E_{\mathbb{P}}[L(S)]\), 
which corresponds to the general robust representation of a coherent risk measure.


Define value at risk and provide a critique of this risk measure in the context of
risk management applications.

Value at Risk at level \(\lambda \in (0,1)\) of a financial position \(X\):
\( \text{V@R}_\lambda(X) := \inf\{m \in \mathbb{R} : \mathbb{P}[X + m < 0] \leq \lambda\} \)

Properties:
- The parameter \(\lambda \in (0,1)\) is chosen close to 0, e.g., \(\lambda = 0.01\) or \(\lambda = 0.005\).
- \(\text{V@R}_\lambda(X)\) is the smallest monetary amount that needs to be added to \(X\) such that the probability 
  of a loss becomes smaller than \(\lambda\) (capital requirement).
- Value at Risk thus allows for a very simple interpretation and can at the same time be easily implemented in practice.

Hint: Value at Risk of the associated loss variable \(L := -X\):

\(
\mathbb{P}[X + m < 0] \leq \lambda \quad \Leftrightarrow \quad F_L(m) := \mathbb{P}[L \leq m] \geq 1 - \lambda
\)

We write (if focussing on losses) 
(ATTENTION: \(\text{V@R}_{1-\lambda}(L) \neq \text{V@R}_{1-\lambda}(X)\)):

\(
\text{V@R}_{1-\lambda}(L) := \inf\{m \in \mathbb{R} : F_L(m) \geq 1 - \lambda\}
\)

V@R neglects extreme events that occur with small probability.
V@R does not generally reward diversification, but charges a larger risk amount for a diversified position in many
cases (no sub-additivity).

In practice, it is a common approach (e.g., for credit risk and underwriting risk) 
to quantify risk in terms of the Mean Value at Risk:

\(\text{MV@R}_\lambda(X) := \text{V@R}_\lambda(X - \mathbb{E}[X]) \quad \text{and} \quad \text{MV@R}_{1-\lambda}(L) := \text{V@R}_{1-\lambda}(L - \mathbb{E}[L])\).

Using cash-invariance this is equivalent to

\(\text{MV@R}_\lambda(X) := \text{V@R}_\lambda(X) - \mathbb{E}[-X] \quad \text{and} \quad \text{MV@R}_{1-\lambda}(L) := \text{V@R}_{1-\lambda}(L) - \mathbb{E}[L]\),

i.e., expected losses are subtracted.
In this case, the capital requirement is only a buffer for ‘unexpected losses’.

The definition of the Solvency Capital Requirement (SCR) in the regulation scheme Solvency II 
for insurance firms is interpreted in the sense of Mean Value at Risk (although the legal framework is self-contradictory).


Give 2 examples highlighting defficiences of V@R 

Which of the following two financial positions is more ‘risky’?
\[
X_1 = \begin{cases} 
1 & \text{with probability } 0.99 \\ 
-1 & \text{with probability } 0.01 
\end{cases}
\]

\[
X_2 = \begin{cases} 
1 & \text{with probability } 0.99 \\ 
-10^{10} & \text{with probability } 0.01 
\end{cases}
\]
The V@R at level 0.01 of both financial positions is -1. However, it is obvious that \(X_2\) is more risky than \(X_1\).

Let \(X_1, X_2\) be two independent financial positions:

\[X_i =
\begin{cases} 
1 & \text{with probability } 0.5 \\ 
-1 & \text{with probability } 0.5 
\end{cases} \quad (i = 1, 2)\]

For both \(X_1\) and \(X_2\), the Value at Risk at level 0.5 amounts to \(-1\).
Consider now the diversified portfolio

\[X_1 + X_2 = 
\begin{cases} 
2 & \text{with probability } 0.25, \\ 
0 & \text{with probability } 0.5, \\ 
-2 & \text{with probability } 0.25 
\end{cases}\]

The Value at Risk at level 0.5 is 0.
Conclusion: In this example, Value at Risk is not sub-additive, i.e.,
\(\text{V@R}_{0.5}(X_1 + X_2) > \text{V@R}_{0.5}(X_1) + \text{V@R}_{0.5}(X_2),\)
and it thus penalizes economically meaningful diversification.


Define convex and coherent risk measures

Definition 4.1.7 (Convex and Coherent Risk Measures)

A monetary risk measure \(\rho\) on \(\mathcal{X}\) is called a convex risk measure 
if it is a convex functional and thus does not penalize diversification.

A convex risk measure is called coherent if it satisfies in addition positive homogeneity, i.e.,
\(
\rho(\lambda X) = \lambda \rho(X) \quad \text{for all } X \in \mathcal{X} \text{ and } \lambda \geq 0.
\)

Remarks:
Each coherent risk measure \(\rho\) is sub-additive, i.e.,

\(
\rho(X + Y) \leq \rho(X) + \rho(Y) \quad \text{for all } X, Y \in \mathcal{X}.
\)

Positive Homogeneity is a scaling property.
- Note that this property might be inappropriate for risk measurement, due to concentration and liquidity risk.
- This criticism motivated the notion of convex risk measures, introduced by Föllmer and Schied (2002).

Define robust representation and elaborate on their construction

Let us denote by \(\mathcal{M}_1 := \mathcal{M}_1(\Omega, \mathcal{F})\) the set of all probability measures on \((\Omega, \mathcal{F})\).
Definition 4.1.8 (Robust Representation)
We say that a convex risk measure \(\rho\) has a robust representation if
\(
\rho(X) = \sup_{Q \in \mathcal{M}_1} \{ \mathbb{E}_Q[-X] - \alpha(Q) \} = \sup_{Q \in \mathcal{Q}_\rho} \{ \mathbb{E}_Q[-X] - \alpha(Q) \}
\)

with some penalty function \(\alpha : \mathcal{M}_1 \to \mathbb{R} \cup \{\infty\}\). Here, the family
\(
\mathcal{Q}_\rho := \{ Q \in \mathcal{M}_1 \mid \alpha(Q) < \infty \}
\)

collects the measures \(Q \in \mathcal{M}_1\) with finite penalty; 
in particular, for \(Q \in \mathcal{Q}_\rho\) one requires that \(\mathbb{E}_Q[-X]\) is well-defined for all \(X \in \mathcal{X}\).

Interpretation:
- The measures in \(\mathcal{Q}_\rho\) can be interpreted as plausible probabilistic models that are taken more or less seriously 
  as specified by the penalty function \(\alpha\).
- The capital requirement \(\rho(X)\) is computed as the worst case of the penalized expected loss \(\mathbb{E}_Q[-X]\), 
  taken over all models \(Q \in \mathcal{Q}_\rho\).

Existence of a robust representation:  
- Under mild conditions, a convex risk measure admits a robust representation, due to the Fenchel-Moreau Theorem.  
- If \(\rho\) admits a robust representation in terms of some penalty function \(\alpha\), 
  then the representation also holds for the minimal penalty function \(\alpha^{\text{min}}\) defined by  
\(
\alpha^{\text{min}}(Q) := \sup_{X \in \mathcal{X}} \{ \mathbb{E}_Q[-X] - \rho(X) \} = \sup_{X \in \mathcal{A}_\rho} \mathbb{E}_Q[-X].
\)

Special case – coherent risk measures:  
- In the coherent case we have \(\alpha^{\text{min}}(Q) \in [0, \infty)\), since the acceptance set \(\mathcal{A}_\rho\) is a convex cone.  
- A coherent risk measure \(\rho\) with robust representation thus takes the form \(\rho(X) = \sup_{Q \in \mathcal{Q}_\rho} \mathbb{E}_Q[-X]\) with  
\(
\mathcal{Q}_\rho = \{ Q \in \mathcal{M}_1 \mid \alpha^{\text{min}}(Q) = 0 \}.
\)

Construction principle:  
- Let \(\alpha : \mathcal{M}_1 \to (-\infty, \infty]\) be a penalty function, and let \(\mathcal{Q}_\rho\) 
  be a subset of probability measures on \((\Omega, \mathcal{F})\).  
- Then a convex risk measure \(\rho\) is given by  
\(
    \rho(X) := \sup_{Q \in \mathcal{Q}_\rho} \{ \mathbb{E}_Q[-X] - \alpha(Q) \}.
\)


Define the risk measures average value at risk. State key properties of this risk
measure.

The question continues with, Prove that it dominates V@R at the same level. What is its robust
representation also state its OCE form?

Average Value at Risk at level \(\lambda \in (0, 1]\) for a financial position \(X\):

\[ \text{AV@R}_\lambda(X) = \frac{1}{\lambda} \int_0^\lambda \text{V@R}_\alpha(X) \, d\alpha \]

Important properties:

1. By construction, AV@R accounts for extreme losses, in contrast to V@R.

2. \(\text{AV@R}\(_\lambda\)\) is a coherent risk measure with robust representation

\(
\text{AV@R}_\lambda(X) = \max_{Q \in \mathcal{Q}_\lambda} \mathbb{E}_Q[-X] 
\)
for \(\mathcal{Q}_\lambda := \left\{ Q \ll \mathbb{P} \mid \frac{dQ}{d\mathbb{P}} \leq \frac{1}{\lambda} \right\} \),
cf. Föllmer and Schied (2016), Theorem 4.52.

3. By monotonicity of \(\alpha \mapsto \text{V@R}_\alpha(X)\),
\(
\text{AV@R}_\lambda(X) \geq \text{V@R}_\lambda(X).
\)

More precisely, AV@R\(_\lambda\) is the smallest convex risk measure which is distribution-based and which dominates 
\(\text{V@R}_\lambda\).

For \(\lambda = 1\), AV@R corresponds to the expected loss \(\mathbb{E}_{\mathbb{P}}[-X]\).

Alternative representation for each \(\lambda\)-quantile \(q_X(\lambda)\) of \(X\):
\(
\text{AV@R}_\lambda(X) = \frac{1}{\lambda} \mathbb{E}[(q_X(\lambda) - X)^+] - q_X(\lambda) = \frac{1}{\lambda} \inf_{l \in \mathbb{R}} \{\mathbb{E}[(l - X)^+] - \lambda l\}
\)

The minimization problem allows to calculate \(\text{AV@R}\) 
without knowledge of \(\text{V@R}\).
Special case:
\(
\text{AV@R}_\lambda(X) = \text{V@R}_\lambda(X) + \frac{1}{\lambda} \mathbb{E}[(-X - \text{V@R}_\lambda(X))^+]
\)

Interpretation: The risk measured in terms of \(\text{AV@R}_\lambda\) consists of \(\text{V@R}_\lambda\) 
and an additional 'excess reserve' which quantifies the expected exceedance of \(\text{V@R}_\lambda\). 
This motivates the alternative name Expected Shortfall.

Tail Value at Risk at level \(\lambda \in (0,1)\):

\(
\text{TV@R}_\lambda(X) := \mathbb{E}_{\mathbb{P}}[-X \mid -X > \text{V@R}_\lambda(X)]
\)

Economic interpretation:
\(\text{TV@R}_\lambda(X)\) is the conditional expected loss given that the loss \(-X\) exceeds the Value at Risk \(\text{V@R}_\lambda(X)\).

Remarks:
- Conditional expectation for events \(A\) with \(\mathbb{P}[A] > 0\):
  \(
  \mathbb{E}_{\mathbb{P}}[X \mid A] = \frac{1}{\mathbb{P}[A]} \mathbb{E}_{\mathbb{P}}[X 1_A]
  \)
- The Tail Value at Risk is in general not a convex risk measure.

Relation between TV@R and AV@R
With \(\gamma_\lambda := \mathbb{P}[-X > \text{V@R}_\lambda(X)]/\lambda\), we obtain

\(
\text{AV@R}_\lambda(X) = \gamma_\lambda \cdot \text{TV@R}_\lambda(X) + (1 - \gamma_\lambda) \cdot \text{V@R}_\lambda(X)
\)

Conclusions:

\(\text{AV@R}_\lambda(X) = \text{TV@R}_\lambda(X)\) for all \(\lambda \in (0,1)\) 
such that \(\mathbb{P}[-X = \text{V@R}_\lambda(X)] = 0\)
For all random variables with continuous distribution, 
the risk measures Average Value at Risk and Tail Value at Risk coincide.
Thus, TV@R is coherent on the set of continuous random variables.

Problem:
For modeling purposes, continuous random variables are typically not sufficient.

Define capacities and Choquet integrals. Can you motivate this notion from facts
that you know from probability (just focus on positive random variables)?

A Choquet capacity is a function \(\mu: \mathcal{F} \rightarrow [0, 1]\) such that

1. \(\mu(\emptyset) = 1 - \mu(\Omega) = 0\)
2. \(A, B \in \mathcal{F}, A \subset B \implies \mu(A) \le \mu(B)\)

(so its missing the \( \sigma \)-additivity of probability measures).
Let \( f \in \mathcal{L}^\infty \),
the Choquet integral of \( f \) with respect to \( \mu \) is:

\( \int f d\mu = \int_{0}^\infty \mu(f > x)dx + \int_{-\infty}^0 \mu(f > x) - 1 dx \)
\(x \ge |f|_\infty \implies \{f > x\} = \emptyset \implies \mu(f > x) = 0 \implies 0 \le \int_0^\infty \mu(f > x) dx \le |f|_\infty\)
\(x \le -|f|_\infty \implies \{f > x\} = \Omega \implies \mu(f > x) = 1 \implies 0 \ge \int_{-\infty}^0 (\mu(f > x) - 1) dx \ge -|f|_\infty\)
1. \(\mu(\emptyset) = 1 - \mu(\Omega) = 0\)
2. \(A, B \in \mathcal{F}, A \subset B \implies \mu(A) \le \mu(B)\)

How are comonotonic measurable functions defined? What are comonotonic risk
measures? How can they be characterized in terms of capacities?

Two measurable functions \( X, Y \) on \((\Omega, \mathcal{F})\) are comonotonic if
\(
(X(\omega) - X(\omega'))(Y(\omega) - Y(\omega')) \geq 0 \quad \forall(\omega, \omega') \in \Omega \times \Omega.
\)

A risk measure \(\rho : \mathcal{X} \to \mathbb{R}\) is comonotonic if
\(
\rho(X + Y) = \rho(X) + \rho(Y)
\)
for comonotonic \(X, Y \in \mathcal{X}\).

Theorem 4.2.5 (Comonotonic Risk Measures)
A monetary risk measure \(\rho : \mathcal{X} \to \mathbb{R}\) is comonotonic, if and only if there exists a capacity \(c\) on \((\Omega, \mathcal{F})\) such that
\(
\rho(X) = \int (-X) \, dc
\)

What are distortion risk measures (DRMs)?

The question continues with,
Under which condition on the distor-
tion function are they coherent? How can they be represented in terms of V@Rs
at different levels?

An increasing function \(\psi : [0, 1] \rightarrow [0, 1]\) with \(\psi(0) = 0\) and \(\psi(1) = 1\) is called a distortion function. 
If \(\mathbb{P}\) is a probability measure on \((\Omega, \mathcal{F})\), then 

\(c^{\psi}(A) := \psi(\mathbb{P}[A]), \quad A \in \mathcal{F}\),
defines a capacity. The risk measure 

\(\rho^{\psi}(X) := \int (-X) \, dc^{\psi}\)
is called a distortion risk measure.

Remark:
If \(\psi\) is left-continuous, then \(\rho^{\psi}(X) = \int (-X) \, dc^{\psi} = \int_{[0,1]} \text{V@R}_{\lambda}(X) \psi'(d\lambda)\). 
For the general case, we refer to Dhaene, Kukush, Linders & Tang (2012).

The distortion risk measure \(\rho^{\psi}(X)\) is coherent, if and only if \(\psi\) is concave.

Provide examples of DRMs together with their distortion functions.

Let \(\alpha, \beta > 0\).
1. Value at Risk: \(\text{V@R}_{\alpha}(X) = \inf\{m \in \mathbb{R} : \mathbb{P}[m + X < 0] \leq \alpha\}\)
2. Average Value at Risk: \(\text{AV@R}_{\beta}(X) = \frac{1}{\beta} \int_{0}^{\beta} \text{V@R}_{\gamma}(X) \, d\gamma\)
3. Range Value at Risk: \(\text{RV@R}_{\alpha, \beta}(X) = \frac{1}{\beta} \int_{\alpha}^{\alpha + \beta} \text{V@R}_{\gamma}(X) \, d\gamma\)

Risk Measure
\(\text{V@R}_{\alpha}\)
\[
    \psi(x) = 
    \begin{cases} 
    0, & 0 \leq x \leq \alpha \\ 
    1, & \alpha < x 
    \end{cases}
\]
\(\text{AV@R}_{\beta}\)
\[
    \psi(x) = 
    \begin{cases} 
    \frac{x}{\beta}, & 0 \leq x \leq \beta \\ 
    1, & \beta < x 
    \end{cases}
\]
\(\text{RV@R}_{\alpha, \beta}\)
\[
    \psi(x) = 
    \begin{cases} 
    0, & 0 \leq x \leq \alpha \\ 
    \frac{x - \alpha}{\beta}, & \alpha < x \leq \alpha + \beta \\ 
    1, & \alpha + \beta < x 
    \end{cases}
\]


Define the MINVAR

For a positive integer \(n\) consider the concave distortion \(c = \psi \circ \mathbb{P}\), 
where \(\psi(x) = 1 - (1 - x)^n\).

The corresponding risk measure \(\rho\), defined as the Choquet integral of the loss, is sometimes called MINVAR.
For independent copies \(X_1, \ldots, X_n\) of \(X\) we get

\(
\rho(X) = \mathbb{E}_{\mathbb{P}}[\max\{-X_1, \ldots, -X_n\}] = -\mathbb{E}_{\mathbb{P}}[\min\{X_1, \ldots, X_n\}],
\)

i.e., \(\rho(X)\) can be described as the expectation under \(\mathbb{P}\) 
of the maximal loss occurring in the portfolio \(X_1, \ldots, X_n\).

How is the risk measure utility-based shortfall risk (UBSR) defined? How does
the coherent special case look like? What form does the acceptance set take in the
coherent case?

Construction:

Let \( l : \mathbb{R} \to \mathbb{R} \) be a convex and increasing loss function.
Let \( r_0 > \inf l \) be a given upper bound for the expected loss \(\mathbb{E}_{\mathbb{P}}[l(-X)]\) of financial positions \( X \in \mathcal{X} \).
Convex acceptance set: \(\mathcal{A} := \{X \in \mathcal{X} | \mathbb{E}_{\mathbb{P}}[l(-X)] \leq r_0 \}\).

Shortfall Risk Measure: 
\(\rho^{SR}(X) = \inf \{m | X + m \in \mathcal{A}\}\)

In terms of the utility function \( u(x) := -l(x) \), the acceptance set can be rewritten as
\(\mathcal{A} = \{X \in \mathcal{X} | \mathbb{E}_{\mathbb{P}}[u(X)] \geq -r_0 \}\).
Therefore, \(\rho^{SR}\) is also called Utility-Based Shortfall Risk (UBSR) measure.

Numerical aspects:
\( m = \rho^{SR}(X) \) is the unique solution of the equation \(\mathbb{E}_{\mathbb{P}}[l(-X - m)] = r_0\).
Stochastic root finding methods (Robbins-Monro algorithm, Polyak-Ruppert algorithm) can be used for the explicit calculation.

Robust representation: (see, e.g., Föllmer and Schied (2016), Proposition 4.113 and Theorem 4.115)
\(
\rho(X) = \max_{Q \in \mathcal{M}_1(\mathbb{P})} \{\mathbb{E}_{Q}[-X] - \alpha^{\min}(Q)\}
\)

with minimal penalty function
\(
\alpha^{\min}(Q) = \inf_{\lambda > 0} \frac{1}{\lambda} \{r_0 + \mathbb{E}_{\mathbb{P}}[l^*(\lambda \frac{dQ}{d\mathbb{P}})]\},
\)
where \( l^*(z) := \sup_{x \in \mathbb{R}} \{zx - l(x)\} \) is the Fenchel-Legendre transform of \( l \).


Define Expectiles and Expectile Value at Risk
For a random variable \(X\), the lower and the upper quantile \(q_{X}(\lambda)_{-}\) and \(q_X(\lambda)_{+}\) 
at level \(\lambda \in (0,1)\) can be defined by the minimization of an asymmetric, piecewise linear loss function in the following sense:

\[
[q_{X}(\lambda)_{-}, q_X(\lambda)_{+}] = \argmin_{x \in \mathbb{R}} \left(\lambda \mathbb{E}[(X-x)_{+}] + (1-\lambda) \mathbb{E}[(X-x)_{-}]\right).
\]

In a similar manner, an expectile is the minimizer of the asymmetric quadratic expected loss
\(
E_\lambda(X) := \argmin_{x \in \mathbb{R}} \left(\lambda \mathbb{E}[(X-x)_{+}]^2 + (1-\lambda) \mathbb{E}[(X-x)_{-}]^2\right).
\)

For \(\lambda=1/2\), we obtain \(E_{1/2}(X) = \mathbb{E}[X]\).
Thus, expectiles can be viewed as asymmetric generalizations of the expectation value.
The word 'expectile' combines 'expectation' and 'quantile'.


Define the entropic risk measure

Convex Entropic Risk Measure with parameter \(\gamma > 0\):

\(e_{\gamma}(X) := \sup_{Q \in \mathcal{M}_1} \left\{\mathbb{E}_{Q}[-X] - \frac{1}{\gamma} H(Q|P)\right\}\)

Construction:
This definition is based on the robust representation of convex risk measures and 
the specification of a penalty function \(\alpha(Q) := \frac{1}{\gamma} H(Q|P)\), defined in terms of the relative entropy

\[H(Q|P) := \begin{cases} 
\mathbb{E}_{Q}[\log \frac{dQ}{dP}] & \text{if } Q \ll P, \\
+\infty & \text{else.}
\end{cases}\]

The relative entropy satisfies \(H(P|P) = 0\), and it is often interpreted as 'distance' between two probability measures. 
However, \(H(\cdot|\cdot)\) is not a metric.

Explicit representation:
The relative entropy satisfies a variational principle (see Föllmer and Schied 2016, Lemma 3.29):

\(H(Q|P) = \sup_{X \in L^{\infty}(\Omega, \mathcal{F}, P)} \left\{\mathbb{E}_{Q}[-X] - \log \mathbb{E}_{P}[e^{-X}]\right\}.\)

This implies \(\mathbb{E}_{Q}[-X] - \frac{1}{\gamma} H(Q|P) \leq \frac{1}{\gamma} \log \mathbb{E}_{P}[e^{-\gamma X}]\), and this upper bound is attained for the measure \(Q\) with density \(e^{-\gamma X}/\mathbb{E}_{P}[e^{-\gamma X}]\).

This yields the explicit representation

\(e_{\gamma}(X) = \frac{1}{\gamma} \log \mathbb{E}_{P}[e^{-\gamma X}]\).


Relationship to the Exponential Premium Principle in actuarial mathematics:  
\(S\) claims amount \(\Rightarrow\) insurance premium 
\(\pi(S) = \frac{1}{\gamma} \log \mathbb{E}_{\mathbb{P}}[e^{\gamma S}] = e_{\gamma}(-S)\)  

Acceptance set:  
\(\mathcal{A} = \{X | e_{\gamma}(X) \leq 0\}\)  
\(= \{X | \mathbb{E}_{\mathbb{P}}[l(-X)] \leq 1\} = \{X | \mathbb{E}_{\mathbb{P}}[u(X)] \geq 0\}\)  

for \(l(x) = e^{\gamma x}\) (loss function) and \(u(x) = 1 - e^{-\gamma x}\) (exponential utility function)  

Conclusion: The Entropic Risk Measure is a special case of the Utility-Based Shortfall Risk.


What is g-divergence? How are divergence
risk measures defined?

Let \( g : \mathbb{R}_+ \to \mathbb{R} \cup \{\infty\} \) be a lower semicontinuous and convex function with \( g(1) < \infty \) 
and superlinear growth \( g(x)/x \to \infty \) as \( x \to \infty \)
For \(\mathbb{Q}, \mathbb{P} \in \mathcal{M}_1\) define the \( g \)-divergence of \(\mathbb{Q}\) w.r.t. \(\mathbb{P}\) by
\[
D_g(\mathbb{Q} | \mathbb{P}) := \begin{cases} 
\mathbb{E}_{\mathbb{Q}} \left[ g \left( \frac{d\mathbb{Q}}{d\mathbb{P}} \right) \right] & \text{if } \mathbb{Q} \ll \mathbb{P}, \\
+\infty & \text{else.}
\end{cases}
\]

Definition 4.2.9 (Divergence Risk Measure)
If \(\mathbb{P}\) is a benchmark model, the resulting convex risk measure \(\rho_g\) defined by
\(
\rho_g(X) = \sup_{\mathbb{Q} \in \mathcal{M}_1(\mathbb{P})} \{\mathbb{E}_{\mathbb{Q}}[-X] - D_g(\mathbb{Q} | \mathbb{P})\}
\)

is called divergence risk measure.
\( g^*(y) = \sup_{x > 0} \{xy - g(x)\} \) convex conjugate of \( g \)
The divergence risk measure can be represented by the variational identity:

\(
\rho_g(X) = \inf_{y \in \mathbb{R}} \{\mathbb{E}_{\mathbb{P}}[g^*(y - X)] - y\}
\)

Provide some prominent examples of divergence risk measures. Why are the divergence
risk measures up to a change of sign equal to the optimized certainty equivalents (OCE's)?

AV@R: 
For \(\lambda \in (0,1]\) let \(g(x) = 0\) for \(x \leq \lambda^{-1}\) and \(g(x) = \infty\) otherwise, 
then \(\rho_g(X) = \text{AV@R}_{\lambda}(X)\).

Entropic risk measure: If \(g(x) = x \log x\), the g-divergence is the relative entropy \(H(\mathbb{Q}|\mathbb{P})\) of \(\mathbb{Q}\) 
w.r.t. \(\mathbb{P}\), and \(\rho_g(X) = \log \mathbb{E}_{\mathbb{P}} [e^{-X}]\).

Optimized Certainty Equivalent:
For a utility function \(u: \mathbb{R} \to \mathbb{R} \cup \{-\infty\}\) with \(u(0) = 0\), 
the optimized certainty equivalent of a financial position \(X \in \mathcal{X}\) is defined as
\(S_u(X) := \sup_{\eta \in \mathbb{R}} \{\eta + \mathbb{E}_{\mathbb{P}} [u(X - \eta)]\).

This can be interpreted as the present value of an optimal allocation of the uncertain future income \(X\) 
between a certain present amount \(\eta\) and the uncertain future amount \(X - \eta\).
Denote by \(g(z) := \sup_{x \in \mathbb{R}} \{xz - l(x)\}\) the convex conjugate function of the loss function \(l\) associated to 
\(u\) via \(l(x) = -u(-x)\). Then the variational identity yields \(S_u(X) = -\rho_g(X)\),
i.e., the optimized certainty equivalent coincides, up to a change of sign, with the divergence risk measure \(\rho_g\).
